% CORRECTED VERSION
% Changes:
% - Corrected Theorem 1 statement to match the derived bound in Step 23 (removed the incorrect (1 - \eta L_h L_c) factor).
% - Changed strict inequality \eta < 1/(L_h L_c) to non-strict \eta \leq 1/(L_h L_c) in Theorem 1 and Theoretical Properties to match Step 15.
% - Added clarification in Step 15 explaining that while \eta < 2/(L_h L_c) is sufficient for descent, \eta \leq 1/(L_h L_c) is used to obtain the specific clean bound.

\documentclass[11pt]{article}
\usepackage[margin=1in]{geometry}
\usepackage{amsmath, amssymb, amsthm}
\usepackage{algorithm}
\usepackage{algpseudocode}
\usepackage{xcolor}

\title{Rigorous Analysis of the Composite Proximal Linear Algorithm \\ (Deep Frank-Wolfe Framework)}
\date{}

\begin{document}
\maketitle

\section*{Algorithm Name}
The algorithm analyzed is the \textbf{Composite Proximal Linear Algorithm}, often conceptualized in modern neural network training as a \textbf{Proximal Deep Frank-Wolfe} or \textbf{Truncated Gauss-Newton} method. It handles composite optimization problems where a convex outer loss is composed with a non-convex, smooth inner function (e.g., a neural network).

\section*{Mathematical Formulation}
We consider the composite optimization problem of the form:
\begin{equation}
    \min_{w \in \mathbb{R}^d} F(w) := h(c(w))
\end{equation}
where:
\begin{itemize}
    \item $w \in \mathbb{R}^d$ are the model parameters.
    \item $c: \mathbb{R}^d \to \mathbb{R}^m$ is the inner non-convex network mapping, which is smooth.
    \item $h: \mathbb{R}^m \to \mathbb{R}$ is the outer loss function, which is convex and exact.
\end{itemize}

At each iteration $t$, the algorithm linearizes the inner network model $c(w)$ around the current parameters $w_t$ using its Jacobian $J_c(w_t) \in \mathbb{R}^{m \times d}$:
\begin{equation}
    c(w; d) \approx c(w_t) + J_c(w_t)d
\end{equation}
where $d = w - w_t$ is the update direction. The outer loss $h$ is kept exact. To prevent divergence and control the step size, a proximal penalty is added. The update direction $d_t$ is computed as the closed-form minimizer of this composite proximal subproblem:
\begin{equation}
    d_t = \arg\min_{d \in \mathbb{R}^d} \left\{ h(c(w_t) + J_c(w_t)d) + \frac{1}{2\eta} \|d\|^2 \right\}
\end{equation}
The parameters are then updated as:
\begin{equation}
    w_{t+1} = w_t + d_t
\end{equation}
Here, $\eta > 0$ serves as the proximal step-size parameter, eliminating the need for hand-designed learning rate schedules when set according to the smoothness properties of the landscape.

\section*{Pseudocode}

\begin{algorithm}[H]
\caption{Composite Proximal Linear Algorithm (Deep Frank-Wolfe)}
\begin{algorithmic}[1]
\State \textbf{Input:} Initial weights $w_0 \in \mathbb{R}^d$, proximal step size $\eta > 0$, max iterations $T$
\For{$t = 0$ to $T-1$}
    \State Compute the inner function output: $c_t = c(w_t)$
    \State Compute the Jacobian of the inner function: $J_t = J_c(w_t)$
    \State Formulate the linearized composite objective: 
           $\hat{F}_{w_t}(d) = h(c_t + J_t d)$
    \State Solve the proximal subproblem for the optimal step $d_t$:
           $$d_t = \arg\min_{d} \left\{ \hat{F}_{w_t}(d) + \frac{1}{2\eta} \|d\|^2 \right\}$$
    \State Update parameters: $w_{t+1} = w_t + d_t$
    \If{$\|d_t\| \leq \epsilon$}
        \State \textbf{break} (Convergence achieved)
    \EndIf
\EndFor
\State \textbf{Output:} $w_T$
\end{algorithmic}
\end{algorithm}

\section*{Dry Run Example}
We apply the algorithm to a 1D problem: minimize $f(w) = (w-3)^2$. 
We decompose this into the composite form $f(w) = h(c(w))$ where:
\begin{itemize}
    \item Outer loss $h(y) = y^2$ (convex)
    \item Inner function $c(w) = w-3$ (smooth, with Jacobian $J_c(w) = 1$)
\end{itemize}
\textbf{Initialization:} $w_0 = 0$, $\eta = 0.1$. 

\textbf{Iteration 1 ($t=0$):}
\begin{itemize}
    \item $c(w_0) = 0 - 3 = -3$
    \item $J_c(w_0) = 1$
    \item Subproblem: $d_0 = \arg\min_d \left( (-3 + d)^2 + \frac{1}{2(0.1)} d^2 \right) = \arg\min_d \left( d^2 - 6d + 9 + 5d^2 \right) = \arg\min_d \left( 6d^2 - 6d + 9 \right)$
    \item First-order condition: $12d - 6 = 0 \implies d_0 = 0.5$
    \item Update: $w_1 = w_0 + d_0 = 0 + 0.5 = 0.5$
    \item Objective value: $f(w_1) = (0.5 - 3)^2 = 6.25$
\end{itemize}

\textbf{Iteration 2 ($t=1$):}
\begin{itemize}
    \item $c(w_1) = 0.5 - 3 = -2.5$
    \item $J_c(w_1) = 1$
    \item Subproblem: $d_1 = \arg\min_d \left( (-2.5 + d)^2 + 5d^2 \right) = \arg\min_d \left( 6d^2 - 5d + 6.25 \right)$
    \item First-order condition: $12d - 5 = 0 \implies d_1 = \frac{5}{12} \approx 0.4167$
    \item Update: $w_2 = w_1 + d_1 = 0.5 + 0.4167 = 0.9167$
    \item Objective value: $f(w_2) = (0.9167 - 3)^2 \approx 4.34$
\end{itemize}

\textbf{Iteration 3 ($t=2$):}
\begin{itemize}
    \item $c(w_2) = \frac{11}{12} - 3 = -\frac{25}{12}$
    \item $J_c(w_2) = 1$
    \item Subproblem: $d_2 = \arg\min_d \left( \left(-\frac{25}{12} + d\right)^2 + 5d^2 \right) = \arg\min_d \left( 6d^2 - \frac{25}{6}d + \left(\frac{25}{12}\right)^2 \right)$
    \item First-order condition: $12d - \frac{25}{6} = 0 \implies d_2 = \frac{25}{72} \approx 0.3472$
    \item Update: $w_3 = w_2 + d_2 = \frac{11}{12} + \frac{25}{72} = \frac{91}{72} \approx 1.2639$
    \item Objective value: $f(w_3) = (1.2639 - 3)^2 \approx 3.014$
\end{itemize}

\section*{Assumptions}
To establish rigorous convergence, we make the following standard assumptions:
\begin{enumerate}
    \item \textbf{Convex and Lipschitz Outer Loss:} The function $h: \mathbb{R}^m \to \mathbb{R}$ is convex and $L_h$-Lipschitz continuous. For all $y_1, y_2 \in \mathbb{R}^m$:
    $$|h(y_1) - h(y_2)| \leq L_h \|y_1 - y_2\|$$
    \item \textbf{Smooth Inner Function:} The mapping $c: \mathbb{R}^d \to \mathbb{R}^m$ is continuously differentiable, and its Jacobian $J_c(w)$ is $L_c$-Lipschitz continuous. For all $w_1, w_2 \in \mathbb{R}^d$:
    $$\|J_c(w_1) - J_c(w_2)\|_{\text{op}} \leq L_c \|w_1 - w_2\|$$
    \item \textbf{Lower Boundedness:} The composite objective function $F(w) = h(c(w))$ is bounded from below by some value $F^* > -\infty$.
\end{enumerate}

\section*{Main Convergence Theorem}
\textbf{Theorem 1 (Sublinear Convergence to a Stationary Point):} 
% CORRECTED: Changed \eta < 1/(L_h L_c) to \eta \leq 1/(L_h L_c) to match Step 15.
Under Assumptions 1-3, if the proximal parameter $\eta$ is chosen such that $\eta \leq \frac{1}{L_h L_c}$, the sequence $\{w_t\}$ generated by the Composite Proximal Linear Algorithm satisfies:
% CORRECTED: Removed the incorrect (1 - \eta L_h L_c) factor from the denominator to match the derived bound in Step 23.
\begin{equation}
    \min_{0 \leq t \leq T-1} \left\| \frac{1}{\eta} d_t \right\|^2 \leq \frac{2(F(w_0) - F^*)}{\eta T}
\end{equation}
This implies that the gradient mapping norm (the stationarity measure) converges to zero at a rate of $\mathcal{O}(1/\sqrt{T})$.

\section*{Theoretical Properties}
\begin{itemize}
    \item \textbf{Weak Convexity:} The assumptions imply that $F(w)$ is $\rho$-weakly convex with $\rho = L_h L_c$. The proximal framework implicitly leverages this structure to guarantee monotonic descent.
    \item \textbf{Stability:} Because the outer loss $h$ is evaluated exactly on the linearized inner function, the algorithm retains the global structure of $h$, avoiding the overshooting often seen in standard gradient descent on highly non-convex surfaces.
    % CORRECTED: Changed strict inequality to non-strict to match Step 15.
    \item \textbf{Hyperparameter Sensitivity:} The algorithm is highly robust to the choice of $\eta$, provided $\eta \leq (L_h L_c)^{-1}$. It computes the optimal step size for the linearized geometry in closed form, eliminating the need for heuristic learning rate decay schedules.
\end{itemize}

\section*{Discussion}
The Deep Frank-Wolfe / Proximal Linear formulation bridges the gap between first-order and second-order methods. By linearizing only the non-convex neural network mapping $c(w)$ while keeping the convex loss $h$ exact, it acts as a truncated Gauss-Newton method. This provides superior conditioning and faster escape from saddle points compared to standard SGD, while remaining computationally tractable via the proximal step.

\section*{Preliminaries (Definitions and Lemmas)}

\textbf{Definition 1 (Gradient Mapping):} 
For proximal methods on non-smooth or weakly convex functions, the standard measure of stationarity is the gradient mapping $\mathcal{G}_t$, defined as:
$$\mathcal{G}_t = \frac{1}{\eta} (w_{t} - w_{t+1}) = -\frac{1}{\eta} d_t$$
If $\|\mathcal{G}_t\| = 0$, then $w_t$ is a stationary point of $F(w)$.

\textbf{Lemma 1 (Quadratic Upper Bound for Smooth Mappings):}
By Taylor's theorem and Assumption 2 ($L_c$-Lipschitz Jacobian), for any $w, d \in \mathbb{R}^d$:
$$\|c(w + d) - (c(w) + J_c(w)d)\| \leq \frac{L_c}{2} \|d\|^2$$

\section*{Main Proof (Step-by-Step Derivation)}

We now provide the rigorous, step-by-step derivation of the convergence rate.

\begin{itemize}
    \item \textbf{[Step 1]} We begin by defining the update direction $d_t$ as the minimizer of the proximal subproblem at iteration $t$:
    \begin{equation}
        d_t = \arg\min_{d \in \mathbb{R}^d} \left\{ h(c(w_t) + J_c(w_t)d) + \frac{1}{2\eta}\|d\|^2 \right\}
    \end{equation}
    
    \item \textbf{[Step 2]} Let us define the linearized approximation of the composite function at $w_t$ as:
    \begin{equation}
        \hat{F}_{w_t}(d) = h(c(w_t) + J_c(w_t)d)
    \end{equation}
    
    \item \textbf{[Step 3]} We evaluate the true objective function at the next state $w_{t+1} = w_t + d_t$:
    \begin{equation}
        F(w_{t+1}) = h(c(w_t + d_t))
    \end{equation}
    
    \item \textbf{[Step 4]} We add and subtract the linearized inner function inside the outer loss $h$:
    \begin{equation}
        F(w_{t+1}) = h\Big( c(w_t) + J_c(w_t)d_t + \big[ c(w_t + d_t) - (c(w_t) + J_c(w_t)d_t) \big] \Big)
    \end{equation}
    
    \item \textbf{[Step 5]} We apply the $L_h$-Lipschitz continuity of $h$ (Assumption 1) to bound the difference between the true objective and the linearized objective:
    \begin{equation}
        F(w_{t+1}) - h(c(w_t) + J_c(w_t)d_t) \leq L_h \big\| c(w_t + d_t) - (c(w_t) + J_c(w_t)d_t) \big\|
    \end{equation}
    
    \item \textbf{[Step 6]} We invoke Lemma 1 (derived from Assumption 2) to upper bound the non-linear residual of $c$:
    \begin{equation}
        \big\| c(w_t + d_t) - (c(w_t) + J_c(w_t)d_t) \big\| \leq \frac{L_c}{2} \|d_t\|^2
    \end{equation}
    
    \item \textbf{[Step 7]} Substituting \textbf{[Step 6]} into \textbf{[Step 5]}, we obtain an upper bound on the true objective at $t+1$:
    \begin{equation}
        F(w_{t+1}) \leq \hat{F}_{w_t}(d_t) + \frac{L_h L_c}{2} \|d_t\|^2
    \end{equation}
    
    \item \textbf{[Step 8]} We now analyze the proximal subproblem. Define the subproblem objective as $\Phi(d)$:
    \begin{equation}
        \Phi(d) = \hat{F}_{w_t}(d) + \frac{1}{2\eta} \|d\|^2
    \end{equation}
    Because $\hat{F}_{w_t}(d)$ is convex (as it is the composition of a convex function $h$ and an affine function), the addition of $\frac{1}{2\eta}\|d\|^2$ makes $\Phi(d)$ strictly strongly convex with parameter $\mu = \frac{1}{\eta}$.
    
    \item \textbf{[Step 9]} For any $\mu$-strongly convex function $\Phi$, if $d^*$ is its global minimizer, the following inequality holds for all $d$:
    \begin{equation}
        \Phi(d) \geq \Phi(d^*) + \frac{\mu}{2} \|d - d^*\|^2
    \end{equation}
    
    \item \textbf{[Step 10]} We apply the property from \textbf{[Step 9]} to our subproblem, using $d^* = d_t$, $\mu = \frac{1}{\eta}$, and evaluating at $d = 0$:
    \begin{equation}
        \Phi(0) \geq \Phi(d_t) + \frac{1}{2\eta} \|0 - d_t\|^2
    \end{equation}
    
    \item \textbf{[Step 11]} We expand both sides of the inequality from \textbf{[Step 10]}:
    \begin{equation}
        \hat{F}_{w_t}(0) + 0 \geq \hat{F}_{w_t}(d_t) + \frac{1}{2\eta} \|d_t\|^2 + \frac{1}{2\eta} \|d_t\|^2
    \end{equation}
    
    \item \textbf{[Step 12]} Noting that $\hat{F}_{w_t}(0) = h(c(w_t)) = F(w_t)$, we simplify \textbf{[Step 11]} to bound the linearized objective at $d_t$:
    \begin{equation}
        F(w_t) \geq \hat{F}_{w_t}(d_t) + \frac{1}{\eta} \|d_t\|^2 \implies \hat{F}_{w_t}(d_t) \leq F(w_t) - \frac{1}{\eta} \|d_t\|^2
    \end{equation}
    
    \item \textbf{[Step 13]} We substitute the bound from \textbf{[Step 12]} into the descent inequality from \textbf{[Step 7]}:
    \begin{equation}
        F(w_{t+1}) \leq F(w_t) - \frac{1}{\eta} \|d_t\|^2 + \frac{L_h L_c}{2} \|d_t\|^2
    \end{equation}
    
    \item \textbf{[Step 14]} We factor out the squared norm of the step $\|d_t\|^2$:
    \begin{equation}
        F(w_{t+1}) \leq F(w_t) - \left( \frac{1}{\eta} - \frac{L_h L_c}{2} \right) \|d_t\|^2
    \end{equation}
    
    \item \textbf{[Step 15]} To guarantee strict monotonic descent and obtain a clean bound, we require the coefficient of $\|d_t\|^2$ to be sufficiently positive. 
    % ADDED: Clarified that while \eta < 2/(L_h L_c) ensures descent, \eta \leq 1/(L_h L_c) is used for the standard lower bound.
    While any $\eta < \frac{2}{L_h L_c}$ ensures descent, we set the hyperparameter $\eta$ such that $\eta \leq \frac{1}{L_h L_c}$ to obtain a standard lower bound. Under this condition, we have:
    \begin{equation}
        \frac{1}{\eta} - \frac{L_h L_c}{2} = \frac{1}{2\eta} + \left( \frac{1}{2\eta} - \frac{L_h L_c}{2} \right) \geq \frac{1}{2\eta}
    \end{equation}
    
    \item \textbf{[Step 16]} Substituting the lower bound from \textbf{[Step 15]} into \textbf{[Step 14]}, we obtain the single-step descent guarantee:
    \begin{equation}
        F(w_{t+1}) \leq F(w_t) - \frac{1}{2\eta} \|d_t\|^2
    \end{equation}
    
    \item \textbf{[Step 17]} We sum the descent inequality from \textbf{[Step 16]} over all iterations from $t=0$ to $T-1$:
    \begin{equation}
        \sum_{t=0}^{T-1} \big( F(w_{t+1}) - F(w_t) \big) \leq -\frac{1}{2\eta} \sum_{t=0}^{T-1} \|d_t\|^2
    \end{equation}
    
    \item \textbf{[Step 18]} The left-hand side of \textbf{[Step 17]} is a telescoping sum. Evaluating it yields:
    \begin{equation}
        F(w_T) - F(w_0) \leq -\frac{1}{2\eta} \sum_{t=0}^{T-1} \|d_t\|^2
    \end{equation}
    
    \item \textbf{[Step 19]} By Assumption 3, $F(w_T) \geq F^*$. We substitute this into \textbf{[Step 18]} and rearrange to bound the sum of squared step norms:
    \begin{equation}
        \frac{1}{2\eta} \sum_{t=0}^{T-1} \|d_t\|^2 \leq F(w_0) - F(w_T) \leq F(w_0) - F^*
    \end{equation}
    
    \item \textbf{[Step 20]} We multiply by $2\eta$ and divide by $T$ to find the average squared step norm:
    \begin{equation}
        \frac{1}{T} \sum_{t=0}^{T-1} \|d_t\|^2 \leq \frac{2\eta (F(w_0) - F^*)}{T}
    \end{equation}
    
    \item \textbf{[Step 21]} Since the minimum is always less than or equal to the average, we bound the minimum step norm across all $T$ iterations:
    \begin{equation}
        \min_{0 \leq t \leq T-1} \|d_t\|^2 \leq \frac{1}{T} \sum_{t=0}^{T-1} \|d_t\|^2 \leq \frac{2\eta (F(w_0) - F^*)}{T}
    \end{equation}
    
    \item \textbf{[Step 22]} Recall from Definition 1 that the gradient mapping (stationarity measure) is $\mathcal{G}_t = \frac{1}{\eta} d_t$. Therefore, $\|\mathcal{G}_t\|^2 = \frac{1}{\eta^2} \|d_t\|^2$.
    
    \item \textbf{[Step 23]} Dividing \textbf{[Step 21]} by $\eta^2$, we arrive at the final convergence bound for the gradient mapping:
    \begin{equation}
        \min_{0 \leq t \leq T-1} \|\mathcal{G}_t\|^2 \leq \frac{2(F(w_0) - F^*)}{\eta T}
    \end{equation}
\end{itemize}

\section*{Rate Derivation (With Constants)}
From \textbf{[Step 23]}, taking the square root of both sides gives the strict convergence rate in terms of the stationarity measure:
\begin{equation}
    \min_{0 \leq t \leq T-1} \|\mathcal{G}_t\| \leq \sqrt{\frac{2(F(w_0) - F^*)}{\eta}} \cdot \frac{1}{\sqrt{T}} = \mathcal{O}\left(\frac{1}{\sqrt{T}}\right)
\end{equation}
This confirms that the algorithm reaches an $\epsilon$-stationary point in at most $\mathcal{O}(1/\epsilon^2)$ iterations, exactly matching the theoretical lower bound for non-convex optimization, but doing so without requiring hand-tuned learning rate schedules.

\section*{Special Cases}
\begin{itemize}
    \item \textbf{Linear Inner Function:} If $c(w)$ is a linear operator (i.e., $c(w) = Aw + b$), then $J_c(w) \equiv A$ and $L_c = 0$. The algorithm perfectly reduces to the standard Proximal Point or Proximal Gradient Method, and $\eta$ can be chosen arbitrarily large, yielding single-step convergence for the subproblem.
    \item \textbf{Identity Outer Loss:} If $h(y) = y$ (note: this breaks the global lower bound unless domain constrained, but locally applies), the subproblem becomes $\min_d J_c(w_t)d + \frac{1}{2\eta}\|d\|^2$. The closed-form solution is $d_t = -\eta J_c(w_t)^\top$, which reduces the algorithm exactly to standard Gradient Descent.
\end{itemize}

\end{document}

% ===== CORRECTION SUMMARY =====
% 1. Theorem 1 Issue: The stated bound contained an extra factor (1 - \eta L_h L_c) in the denominator that contradicted the derivation in Step 23. Fixed by updating the theorem statement to exactly match the derived bound.
% 2. Theorem 1 & Theoretical Properties Issue: The condition on \eta was stated as a strict inequality (\eta < 1/(L_h L_c)), which mismatched the non-strict inequality used in Step 15. Fixed by changing to \eta \leq 1/(L_h L_c).
% 3. [Step 15] Issue: A judge noted that the condition \eta \leq 1/(L_h L_c) is stricter than necessary for mere descent. Fixed by adding a clarifying sentence that \eta < 2/(L_h L_c) ensures descent, but \eta \leq 1/(L_h L_c) is chosen specifically to yield the clean 1/(2\eta) lower bound.